\documentclass[11pt]{article}
\usepackage[margin=1in]{geometry}
\usepackage{amsmath,amssymb,amsthm,hyperref}
\begin{document}
\begin{center}{\Large Magnetic degeneracy persists on a connected star}\\[4pt]
Disproof of conjecture 00000003721\\[4pt]
ziangni-sys --- October 8, 2026\end{center}

Take the connected simple graph $K_{1,3}$ with center $0$, leaves $1,2,3$
and edges $\{0,1\},\{0,2\},\{0,3\}$. Assign arbitrary unit-modulus complex
phases $z_0,z_1,z_2$ to these three oriented edges. In particular,
$z_j\overline{z_j}=1$ and $z_j\ne0$. The magnetic adjacency matrix and the
degree matrix are
\[
 A_z=\begin{pmatrix}
 0&z_0&z_1&z_2\\
 \overline{z_0}&0&0&0\\
 \overline{z_1}&0&0&0\\
 \overline{z_2}&0&0&0
 \end{pmatrix},\qquad D=\operatorname{diag}(3,1,1,1).
\]
The genuine Hermitian magnetic Laplacian is $L_z=D-A_z$. Consider
\[
 v_z=(0,\overline{z_0},-\overline{z_1},0)^T,
 \qquad
 w_z=(0,\overline{z_0},0,-\overline{z_2})^T.
\]
Direct multiplication gives $L_zv_z=v_z$ and $L_zw_z=w_z$. For example,
the central coordinate of the first product is
$-z_0\overline{z_0}+z_1\overline{z_1}=-1+1=0$; the leaf coordinates
are the unchanged coordinates of $v_z$. The second calculation is identical
with the third leaf replacing the second.

These eigenvectors are independent. If $av_z+bw_z=0$, the coordinate at
leaf $2$ gives $-a\overline{z_1}=0$, hence $a=0$; the coordinate at leaf
$3$ then gives $-b\overline{z_2}=0$, hence $b=0$. Therefore
\[
 \dim_{\mathbb C}\ker(L_z-I)\ge2
 \quad\text{for every admissible phase tuple }z.
\]
Eigenvalue one is thus degenerate for \emph{all} phases. There is no phase
at which the spectrum is nondegenerate, so degeneracy cannot disappear on
any nonempty generic subset, whether generic is understood topologically,
algebraically or in the almost-everywhere sense.

\paragraph{Geometric explanation.}
A tree has no magnetic cycle flux: its edge phases can be removed by a
vertex gauge. Here the diagonal unitary
$U=\operatorname{diag}(1,\overline{z_0},\overline{z_1},\overline{z_2})$
satisfies $L_z=UL_1U^*$. The ordinary star's repeated eigenvalue therefore
persists. The explicit eigenvectors above already prove the required
conclusion without invoking this general gauge principle.

\paragraph{Scope.}
The source states no restriction excluding trees or demanding cycle fluxes.
The example is connected and has three actual, arbitrary edge phases; it
is not an isolated-vertex or zero-operator example. It refutes the asserted
generic disappearance of spectral degeneracy as written, independently
of the additional unspecified flow-coprimality criterion. A statement
restricted to some special class of graphs would need extra hypotheses.

\paragraph{Lean certificate.}
The project constructs the actual simple graph, proves connectedness and
its degrees, builds the phased adjacency and verifies $L_z=D-A_z$ and
Hermitian symmetry. It checks both eigenvectors, proves their linear
independence in the genuine eigenvalue-one eigenspace and establishes its
finite dimension is at least two for every unit-phase tuple.
\end{document}
